\section{Technical Approach}

\subsection{Pipeline}

Let $\mathcal{C}$ be a bank of English activity sentences, each with a CLIP
text embedding $e_c$. For a captured frame $x_i$, SmolVLM-256M runs on the
smartphone in 4-bit form and writes a short action phrase $s_i$ together
with the probability $p_{i,t}$ of each generated token. The confidence of
the phrase is
\begin{equation}
\kappa_i
=
\min_t p_{i,t}.
\end{equation}
The phrase is shown as soon as this local inference finishes. It does not
wait for the server.

The JPEG is uploaded only when the phrase is unsafe:
\begin{equation}
\kappa_i < \tau_{\mathrm{u}}.
\end{equation}
Otherwise the phone sends the sentence alone. CLIP on the server maps that
sentence into the bank,
\begin{equation}
c_{d,i}
=
\arg\max_{c\in\mathcal{C}}
\cos\big(f_T(s_i),e_c\big),
\end{equation}
and this match is the caption that stays on screen.

When the JPEG is uploaded, CLIP also retrieves a caption from the picture,
\begin{equation}
c_i^\star
=
\arg\max_{c\in\mathcal{C}}
\cos\big(f_I(x_i),e_c\big).
\end{equation}
The phone sentence is kept when $c_i^\star=c_{d,i}$. Otherwise the displayed
caption is replaced by $c_i^\star$. A larger generative model, Qwen2-VL, can
take the same place as CLIP: it writes its own phrase, the phrase is mapped
into $\mathcal{C}$ by the same text encoder, and the caption is replaced
when that match differs from $c_{d,i}$.

Only the low-confidence path puts an image on the air:
\[
\mathrm{Smartphone}
\rightarrow
\mathrm{Wi\mbox{-}Fi}
\rightarrow
\mathrm{server~checker}
\rightarrow
\mathrm{Smartphone}.
\]

\begin{algorithm}[t]
\caption{Local sentence and edge check}
\label{alg:verify}
\begin{algorithmic}[1]
\STATE Run 4-bit SmolVLM on frame $x_i$ and display $s_i$.
\STATE $\kappa_i\leftarrow\min_t p_{i,t}$.
\IF{$\kappa_i\geq\tau_{\mathrm{u}}$}
    \STATE Send $s_i$ only and keep its bank match $c_{d,i}$.
\ELSE
    \STATE Predict the edge checking time $\hat T_i^{\mathrm{edge}}$.
    \STATE Compute the uplink budget $B_i^{\mathrm{up}}$.
    \STATE Wait for an upload grant, then send the JPEG.
    \STATE Retrieve $c_i^\star$ with CLIP, or with Qwen2-VL.
    \IF{$c_i^\star=c_{d,i}$}
        \STATE Accept $c_{d,i}$.
    \ELSE
        \STATE Replace the caption with $c_i^\star$.
    \ENDIF
\ENDIF
\end{algorithmic}
\end{algorithm}

\subsection{Latency Model}

The sentence that appears first is local:
\begin{equation}
T_i^{\mathrm{draft}}
=
T_i^{\mathrm{local}},
\end{equation}
where $T_i^{\mathrm{local}}$ is the SmolVLM execution time on the phone.

A low-confidence frame also waits for a checked caption. Its confirmation
delay is
\begin{equation}
T_i^{\mathrm{conf}}
=
T_i^{\mathrm{wait}}
+
T_i^{\mathrm{up}}
+
T_i^{\mathrm{edge}}
+
T_i^{\mathrm{down}},
\end{equation}
where $T_i^{\mathrm{wait}}$ is the time until the scheduler grants the
upload, $T_i^{\mathrm{up}}$ is the JPEG upload time,
$T_i^{\mathrm{edge}}$ is the server checking time, and
$T_i^{\mathrm{down}}$ is the return time of the result. The downlink carries
an acceptance decision or a caption identifier, so the large terms are the
uplink and the checker. The request meets its deadline when
\begin{equation}
T_i^{\mathrm{conf}}
\leq
D_i^{\mathrm{conf}}.
\end{equation}

\subsection{Online Edge-Latency Prediction}

Checker time is not treated as a constant. It changes with the uploaded
image and with the load on the server. The predictor estimates that time
before the JPEG is uploaded, from
\begin{equation}
\mathbf{z}_i
=
\left[
L_i,\;
U_i^{\mathrm{CPU}},\;
T_{i-1}^{\mathrm{edge}}
\right],
\end{equation}
where $L_i$ is the JPEG size in bytes, $U_i^{\mathrm{CPU}}$ is the server
CPU utilization when the request arrives, and $T_{i-1}^{\mathrm{edge}}$ is
the measured checking time of the previous request. The JPEG size is known
once the phone has encoded the frame. It sets the upload time and is also a
cue to image complexity. CPU utilization is read on the server at the
request. The previous checking time carries warm-up, contention, and
model-dependent variation.

A small network $g_\theta(\cdot)$ predicts
\begin{equation}
\hat T_i^{\mathrm{edge}}
=
g_\theta(\mathbf{z}_i).
\end{equation}
It is trained on traces from the server used in the demo. The network stays
small because it supports scheduling, not captioning. Accuracy is the mean
absolute error between predicted and measured checking time.

\subsection{Compute-Aware Link Control}

Three controls sit above the Wi-Fi MAC. The upload threshold
$\tau_{\mathrm{u}}$ decides whether a picture follows the sentence. JPEG
quality and resolution decide how large that picture is. The scheduler
decides which phone is checked first.

For an upload request $i$, the remaining uplink budget is
\begin{equation}
B_i^{\mathrm{up}}
=
\left[
D_i^{\mathrm{conf}}
-
\hat T_i^{\mathrm{edge}}
-
\hat T_i^{\mathrm{down}}
-
M
\right]^+,
\end{equation}
where $\hat T_i^{\mathrm{down}}$ is the expected return time and $M$ is a
margin for prediction error. If the request reaches the scheduler at time
$a_i$, its uplink completion deadline is
\begin{equation}
d_i^{\mathrm{up}}
=
a_i+B_i^{\mathrm{up}}.
\end{equation}
The scheduler serves the pending request with the earliest deadline:
\begin{equation}
i^\star
=
\arg\min_{i\in\mathcal{Q}}
d_i^{\mathrm{up}}.
\end{equation}
Each phone first sends a short request that contains the JPEG size and the
confirmation deadline, then waits for an application-layer grant. The
granted phone uploads the JPEG. A long predicted edge time leaves less room
for the upload, so that request gets an earlier uplink deadline. A short
predicted edge time can wait longer without missing
$D_i^{\mathrm{conf}}$.